% ===========================================================================
\chapter{Checks, and what next}\label{ch:next}

\section{Checking everything}

Contracts check each run of the engine as it happens. A set of checks that
runs the engine on known cases, and compares its answers with the right
ones, catches what contracts cannot: a parser that reads a program wrongly
but consistently, a rule written correctly that means something other than
intended. \code{checks.nex} runs every part of the engine on small programs
of its own.

\begin{newfile}{checks.nex}
-- checks.nex: executable checks for every piece of the project.
-- Run: nex checks.nex

intern Parser
intern Analysis
intern Query
intern Explain

class Checker
feature
  passed: Integer
  failed: Integer

  check(label: String, ok: Boolean) do
    if ok then
      passed := passed + 1
      console.print_line("PASS " + label)
    else
      failed := failed + 1
      console.print_line("FAIL " + label)
    end
  end

  summary do
    console.print_line("---")
    console.print_line(passed.to_string + " passed, " + failed.to_string + " failed")
    if failed > 0 then
      exit(1)
    end
  end
end
\end{newfile}

A few helpers turn text into the things the checks need:

\begin{continued}{checks.nex}
-- The program `text` describes; it must be one.
function program_of(text: String): Datalog_Program do
  match parse_program(text) of
    Ok as ok then result := ok.value
    Err as e then raise e.error
  end
end

-- The one atom in `text`.
function atom_of(text: String): Atom do
  result := program_of("?- " + text + ".").queries.get(0)
end

-- What parsing `text` reports, or "" if it parses.
function parse_error(text: String): String do
  match parse_program(text) of
    Ok as ok then result := ""
    Err as e then result := e.error
  end
end

-- Every fact `text` implies, as sorted text, evaluated with `e` into `db`.
function evaluated(text: String, db: Database, e: Evaluator): Array[String] do
  evaluate_program(program_of(text), db, e)
  result := []
  across db.predicates as p do
    across db.facts_of(p) as f do
      result.add(f.text)
    end
  end
  result := result.sort
end
\end{continued}

\begin{continued}{checks.nex}
function check_terms(c: Checker) do
  let a := atom_of("parent(X, \"Ann Lee\", 3, X)")
  c.check("an atom prints as it is written", a.text = "parent(X, \"Ann Lee\", 3, X)")
  c.check("its variables, each once", a.variables = ["X"])
  c.check("an atom with a variable is not ground", not a.is_ground)
  c.check("a fact is ground", atom_of("parent(ann, bo)").is_ground)
end

function check_parser(c: Checker) do
  let p := program_of("% a comment
p(a). p(b).
q(X) :- p(X), X != a.
r(X) :- p(X), not q(X).
s(X) :- p(X), t(_, _).
?- q(Who).")
  c.check("facts, rules and questions are read", p.facts.length = 2 and p.rules.length = 3
    and p.queries.length = 1)
  c.check("a rule prints back as written", p.rules.get(1).text = "r(X) :- p(X), not q(X).")
  c.check("each _ is a variable of its own",
    literal_text(p.rules.get(2).body.get(1)) = "t(_1, _2)")
  c.check("errors give the line", parse_error("p(a).\nq(b) :- p(a)\n").starts_with("line "))
  c.check("a fact may not have variables", parse_error("p(X).").contains("cannot have variables"))
  c.check("a stray character is reported", parse_error("p(a) & q.").contains("unexpected character"))
end

function check_matching(c: Checker) do
  let none := create Bindings.empty
  let b := match_atom(atom_of("parent(X, Y)"), atom_of("parent(ann, bo)"), none)
  c.check("a pattern matches a fact", b /= nil and b.text = "X = ann, Y = bo")
  c.check("a repeated variable must match the same value",
    match_atom(atom_of("parent(X, X)"), atom_of("parent(ann, bo)"), none) = nil)
  c.check("constants must be equal",
    match_atom(atom_of("parent(cy, Y)"), atom_of("parent(ann, bo)"), none) = nil)
  c.check("comparisons on integers", holds("<", create Int.make(2), create Int.make(10)))
  c.check("ordering needs integers", not holds("<", create Symbol.make("a"), create Symbol.make("b")))
end
\end{continued}

The evaluation checks compare the three ways of evaluating a program, naive,
semi-naive, and semi-naive with an index, and insist that they derive the
same facts. Contracts already check each one against the fixpoint; these
checks also make sure they agree with one another, and with the answers we
expect.

\begin{continued}{checks.nex}
function check_evaluation(c: Checker) do
  let family := "parent(a, b). parent(b, c). parent(c, d).
ancestor(X, Y) :- parent(X, Y).
ancestor(X, Z) :- parent(X, Y), ancestor(Y, Z)."
  let naive := evaluated(family, create Database.make, create Naive_Evaluator.make)
  let semi := evaluated(family, create Database.make, create Semi_Naive_Evaluator.make)
  let indexed := evaluated(family, create Indexed_Database.make, create Semi_Naive_Evaluator.make)
  c.check("naive evaluation finds every ancestor", naive.contains("ancestor(a, d)") and naive.length = 9)
  c.check("semi-naive evaluation finds the same facts", semi = naive)
  c.check("so does an indexed database", indexed = naive)
  let chain := "link(t1, t2). link(t2, t3). link(t3, t4). link(t4, t5). link(t5, t6).
route(X, Y) :- link(X, Y).
route(X, Z) :- route(X, Y), link(Y, Z)."
  let slow := create Naive_Evaluator.make
  let fast := create Semi_Naive_Evaluator.make
  evaluated(chain, create Database.make, slow)
  evaluated(chain, create Database.make, fast)
  c.check("semi-naive evaluation looks at fewer facts", fast.attempts < slow.attempts)
end

function check_analysis(c: Checker) do
  c.check("a head variable must be bound",
    program_problem(program_of("p(X) :- q(Y).")).contains("variable X in the head"))
  c.check("a test's variables must be bound before it",
    program_problem(program_of("p(X) :- X != a, q(X).")).contains("variable X in 'X != a'"))
  c.check("a negated literal's variables must be bound",
    program_problem(program_of("p(X) :- q(X), not r(Y).")).contains("variable Y"))
  c.check("a cycle through negation is refused",
    program_problem(program_of("q(a). p(X) :- q(X), not p(X).")).contains("its own negation"))
  let access := evaluated("user(ann). user(bo). banned(bo).
ok(U) :- user(U), not banned(U).
refused(U) :- user(U), not ok(U).", create Database.make, create Semi_Naive_Evaluator.make)
  c.check("negation, one stratum at a time", access.contains("ok(ann)")
    and not access.contains("ok(bo)") and access.contains("refused(bo)"))
end

function check_queries(c: Checker) do
  let db := create Database.make
  evaluate_program(program_of("parent(a, b). parent(a, c).
sibling(X, Y) :- parent(P, X), parent(P, Y), X != Y."), db, create Semi_Naive_Evaluator.make)
  c.check("a question with a variable", answer_lines(atom_of("sibling(b, Who)"), db) = ["Who = c"])
  c.check("a question without one", answer_lines(atom_of("sibling(b, c)"), db) = ["yes"]
    and answer_lines(atom_of("sibling(b, b)"), db) = ["no"])
end

function check_explanations(c: Checker) do
  let db := create Explained_Database.make
  evaluate_program(program_of("parent(a, b). parent(b, c).
ancestor(X, Y) :- parent(X, Y).
ancestor(X, Z) :- parent(X, Y), ancestor(Y, Z)."), db, create Semi_Naive_Evaluator.make)
  let proof := db.why(atom_of("ancestor(a, c)"))
  c.check("a derived fact is explained by its rule", proof.get(0).starts_with("ancestor(a, c)    by "))
  c.check("down to the given facts", proof.contains("  parent(a, b)    (given)"))
  c.check("every step checks", db.steps_hold(atom_of("ancestor(a, c)")))
end
\end{continued}

\begin{continued}{checks.nex}
function main() do
  let c := create Checker
  check_terms(c)
  check_parser(c)
  check_matching(c)
  check_evaluation(c)
  check_analysis(c)
  check_queries(c)
  check_explanations(c)
  c.summary
end

main()
\end{continued}

\runoutput{checks.nex}{checks}

\section{What next}

The engine is a complete Datalog, small but real. Here is where the
subject goes from here, and how the engine could follow.

\begin{description}[style=nextline, leftmargin=1.2em, itemsep=0.4em]
  \item[Answering only what is asked.] Our engine derives every fact a
    program implies, then answers questions from them. A question such as
    \code{?- ancestor(alice, X).} only needs Alice's descendants. The
    \emph{magic sets} transformation rewrites a program so that bottom-up
    evaluation derives only the facts relevant to a question, keeping the
    simple evaluator and its guarantees.
  \item[Aggregates.] Counting, summing and finding a minimum (``how many
    hops is the shortest route?'') need care, for the same reason negation
    did: an aggregate over a predicate is only right once that predicate is
    complete. Stratification extends to them.
  \item[Better joins.] Our join tries the body's atoms in the order they
    are written. A real engine reorders them, starting with the atom that
    has the fewest candidates, and indexes on whichever arguments a join
    binds, not just the first.
  \item[Datalog in use.] Soufflé compiles Datalog to fast C++ and is used to
    analyse large programs for bugs and security holes. Datomic and XTDB use
    Datalog as a query language for databases. Reading about either is a
    good way to see how far the ideas in this tutorial go.
\end{description}

Each of these is a natural next step for the engine, and each comes with
the same question this tutorial kept asking: what must be true, and how can
the program check it?
